\chapter{Conclusion}
\chaptershorttitle{Conclusion}
\label{cha:conclusion}

The main insights of this investigation are\ldots
\begin{infobox}[title={Conclusion vs. summary}]
	In the final chapter of the manuscript, you should provide a conclusion of the work. However, in practice many authors formulate a summary instead of a genuine conclusion. Hence, please note the differences:
	\begin{itemize}
		\item \textbf{Conclusion:} Provides a final interpretation of the study's results and offers insights into the broader implications of the findings. It often discusses the significance of the research in the context of existing knowledge. If relevant, key quantitative figures should be included. In simple terms: what are the take-home insights of the work?
		\item \textbf{Summary:} Consecutively recaps all key points of the entire work. It provides a step-by-step overview of the main objectives, assumptions, applied methods, empirical experiments, and results. The summary therefore goes beyond the conclusion, as it focuses not only on the findings but also provides a broader overview of the entire work.
	\end{itemize}
	In this template, the manuscript already starts with an abstract, which is a concise summary. Therefore, the final conclusion should not be a second, redundant summary.
\end{infobox}

% Examples: correct vs. incorrect conclusion
\begin{successbox}[title={Correct conclusion -- example}]
	This thesis shows that the proposed approach improves robustness under realistic noise while meeting real-time constraints. The key takeaway is that combining temporal context with lightweight uncertainty calibration yields consistent gains up to \SI{5.2}{\percent} across diverse datasets. This implies practical deployability on embedded hardware with reduced maintenance effort. Remaining limitations include sensitivity to extreme domain shifts by \SI{42}{\percent}, motivating follow-up work on adaptive calibration.
\end{successbox}

\begin{errorbox}[title={Incorrect conclusion (summary) -- example}]
	In Chapter~1 we introduced the problem, in Chapter~2 we reviewed related work, in Chapter~3 we described the method, in Chapter~4 we presented experiments, and in Chapter~5 we discussed the results. The dataset and training details were also explained. This restates the thesis structure but does not synthesize what the results mean or why they matter; it is a summary, not a conclusion.
\end{errorbox}

Another hint: to ensure that the conclusion is concise and to the point, consider using a bulleted list and avoiding long, complex sentences. This will help the reader extract the main insights quickly. In most cases, a conclusion should not exceed one page.

\section{Future work}
\label{sec:futurework}
The outlook of this work should provide a brief overview of potential future research directions. This can include open questions, unsolved problems, or potential improvements to the current work. It is also possible to discuss potential applications of the findings in practice. A bulleted list is again a good choice for this section to ensure a clear and concise presentation of future work.
